\section{Problem Statement and Objectives}
\label{sec:problem}

\subsection{Thematic domain and motivation}

This project is situated in the financial-economics theme, specifically
at the intersection of Indian monetary policy operations and applied
machine learning on open public data. The Reserve Bank of India
publishes a rich set of weekly statistical tables through the National
Data and Analytics Platform (NDAP), covering the policy rate corridor,
the central-bank balance sheet, and the short-term debt markets. These
are the inputs that determine overnight rupee liquidity, and their
combined behaviour shapes the cost of every short-tenor borrowing in
the formal banking system.

The Weighted Average Call Money Rate (WACMR) is the daily volume-weighted
mean rate at which scheduled commercial banks borrow and lend overnight
to one another in the call-money market. It is the operational target
of the Liquidity Adjustment Facility (LAF), and any sustained deviation
between WACMR and the policy Repo Rate is a real-time diagnostic of
whether the corridor is binding. WACMR moves first when liquidity
conditions tighten or loosen, before the broader rate term structure
follows.

\subsection{Why WACMR matters}

WACMR sits at the head of the entire monetary transmission chain and
carries direct, measurable consequences for three groups.

\begin{itemize}
\item \textbf{Reserve Bank of India.} The WACMR-Repo spread is the
  primary signal of whether the announced policy stance is being
  enforced in the market. Persistent positive spreads indicate
  liquidity deficit and presage interventions through Variable Rate
  Repo (VRR) auctions; persistent negative spreads indicate surplus
  conditions and presage Variable Rate Reverse Repo (VRRR) operations.
\item \textbf{Banks, mutual funds, and corporate treasuries.} Overnight
  fund net asset values, money-market mutual fund returns, and the
  duration risk on short-end bond portfolios all move with WACMR. A
  forecast accurate to within roughly ten basis points is directly
  actionable for intraday positioning and treasury hedging.
\item \textbf{Macroeconomic researchers and policy analysts.}
  Identifying the structural regime boundaries in India's liquidity
  framework, in particular the pre-COVID and post-COVID periods,
  provides quantitative evidence on how exogenous shocks shift the
  monetary-transmission mechanism.
\end{itemize}

\subsection{The forecasting problem}

WACMR is difficult to forecast with simple rules because it is the
equilibrium outcome of five interacting channels operating on different
frequencies. The five channels are the policy corridor (Repo, Reverse
Repo, MSF), system liquidity (RBI balance sheet items, M3, reserve
money), short-term instrument flows (Treasury Bills, G-Secs, Commercial
Paper, Market Repo), the macro price environment (CPI indices), and
exogenous market sentiment (equity and currency). The relative weights
of these channels also change across structural regimes, which makes
the problem non-stationary.

\newpage

\subsection{Hypotheses}

Four hypotheses are tested in this report. Each is stated in advance,
with the supervised and unsupervised stages designed to provide
falsifiable evidence on each.

\begin{enumerate}
\item \textbf{H1 (regime structure).} India's weekly money-market data
  contains at least two statistically distinct macroeconomic regimes,
  identifiable by unsupervised clustering on RBI balance-sheet and
  rate-corridor variables.
\item \textbf{H2 (regime-aware forecasting).} Injecting regime labels
  and centroid distances as features into an XGBoost regressor reduces
  out-of-sample RMSE by at least fifteen percent relative to a baseline
  model with no regime features.
\item \textbf{H3 (corridor dominance).} The Repo, Reverse Repo, and MSF
  rates are the dominant predictors in the supervised model, identified
  by SHAP feature importance.
\item \textbf{H4 (market sentiment marginal impact).} Adding twenty-eight
  Nifty 50 and USD/INR features (OHLCV plus five custom technical
  indicators per instrument) materially improves WACMR forecasting at
  the weekly frequency.
\end{enumerate}

\subsection{Project objectives}

The five objectives below map one-to-one onto the five guidelines parts
of the project specification.

\begin{enumerate}
\item Build a clean, density-filtered, Friday-aligned weekly macro
  panel from the NDAP API and supplemental market data, persisted in a
  queryable SQLite database.
\item Discover statistically distinct monetary regimes from the data
  alone using PCA and K-Means, and validate the result against known
  macroeconomic events.
\item Forecast WACMR one week ahead using walk-forward XGBoost
  regression, with SHAP interpretability, and quantify whether
  regime-awareness improves forecasting accuracy.
\item Translate the technical findings into concrete, evidence-grounded
  policy and trading recommendations.
\item Ship the analysis as a live, interactive system that the user can
  dogfood, including a policy counterfactual simulator and a
  Gemini-powered data agent that can answer free-form questions over
  the dataset.
\end{enumerate}

The remaining sections of this report follow this structure. Section
\ref{sec:data} documents the ten datasets and their sourcing. Section
\ref{sec:eda} presents the exploratory data analysis. Section
\ref{sec:database} details the database design. Section
\ref{sec:methods} sets out the methods, including the regime-discovery
algorithm, the walk-forward forecasting protocol, and the NLP event
corpus. Section \ref{sec:results} interprets the findings as evidence
about the current state of India's overnight liquidity. Section
\ref{sec:recommendations} translates these findings into actionable
recommendations. Sections \ref{sec:system} and \ref{sec:qa} cover the
bonus interactive system, deployment, and quality assurance. Sections
\ref{sec:limitations} and \ref{sec:conclusion} close out the limitations
and the headline takeaway.
